\section{Introduction}\label{intro:section}

The Smith--Toda realization problem asks how far one can kill the chromatic
parameters in Brown--Peterson homology, one at a time, while retaining a finite
spectrum. At a prime $p$, write
\[
 BP_* = \mathbb Z_{(p)}[v_1,v_2,\ldots],
 \qquad |v_i|=2(p^i-1),
\]
with Hazewinkel generators. A \emph{finite $p$-local spectrum} is a retract
of the $p$-localization of a finite CW spectrum. A \emph{Smith--Toda
complex $V(n)$} is such a spectrum with Brown--Peterson homology
$BP_*/(p,v_1,\ldots,v_n)$, including the canonical $BP_*BP$-comodule
structure and a generator in degree zero. In this quotient each listed
generator is killed to its first power. The comodule records the
Brown--Peterson cooperations, so its specification is part of the
realization problem, in addition to the underlying graded module.
The periodicity theorem gives generalized Moore spectra whose homology
kills suitable powers of the initial generators
\citep[Proposition~5.14]{HopkinsSmith1998}. Requiring every exponent
to be one is a more restrictive realization problem.

The mod-$p$ Moore spectrum gives $V(0)$ at every prime. The classical
constructions of $V(1)$, $V(2)$, and $V(3)$ for $p>2$, $p>4$, and $p>6$
are due to Adams, Smith, and Toda
\citep{Adams1966,Smith1970,Toda1971}; see also the historical account in
\citet[p.~493]{Nave2010}. Successive cofibers turn a suitable first-power
$v_i$ self-map into the next quotient. The difficulty is producing the
self-map on the finite spectrum: an algebraically primitive coefficient
need not survive the Adams--Novikov spectral sequence (ANSS) to an actual map.
These constructions thus connect the algebra of complex cobordism with
specific attaching maps in the stable homotopy category.

The next case remained outside those constructions. The September 2023
version of \citet[Remark~3.29]{CulverZhang2024} records the existence of
$V(4)$ as open at every prime. There are also height-dependent
nonexistence results: \citet[Theorem~1.3]{Nave2010} proves that
$V((p+1)/2)$ does not exist for $p\geq7$, and that $V((p-1)/2)$,
if it exists, cannot be a ring spectrum. Those obstructions leave
room for $V(4)$ at a sufficiently large prime.
We establish the following specified-prime realization.

\begin{theorem}\label{intro:main}
At $p=1009$ there is a finite $p$-local spectrum $X$ and an isomorphism
\[
 BP_*X\ \cong\ BP_*/(p,v_1,v_2,v_3,v_4)
\]
of graded $BP_*BP$-comodules, with the canonical quotient comodule on the
right and its generator in degree zero.
\end{theorem}

Theorem~\ref{intro:main} gives a positive answer to the Smith--Toda $V(4)$
existence problem at the stated prime. Its conclusion concerns the ordinary
$p$-local stable category; finiteness is established before any chromatic
localization. For comparison,
\citet[Theorem~1.7]{KatoShimomuraShimomura2025} construct at this prime
an $L_5$-local spectrum $W_5$ with
\[
 BP_*W_5=v_5^{-1}BP_*/(p,v_1,v_2,v_3,v_4),
\]
where $L_5$ denotes localization with respect to $v_5^{-1}BP$.
That realization has $v_5$ inverted and does not supply the ordinary
finite realization of the unlocalized quotient required here.
The choice $1009$ makes the finite inequalities in our calculation
generous. We do not require a product on the final spectrum $X$.

\subsection{Methods and proof strategy}

The algebraic setting comes from the spectrum of Brown and Peterson
\citep{BrownPeterson1966}, Quillen's formal-group description of
cobordism and its $p$-typical summand \citep{Quillen1969}, and
Hazewinkel's explicit $p$-typical generators \citep{HazewinkelIII}.
Our main tool is the sphere-source Adams--Novikov spectral sequence,
whose cobordism-theoretic construction goes back to Novikov
\citep{Novikov1967}. We use its ordinary $BP$ resolution, its cobar
$E_2$ term, and its compatibility with actual binary maps; see
\citet[Sections~2.2--2.3]{Ravenel2004}. The algebraic estimates use
filtered Hopf-algebra cohomology as developed by
\citet[Section~4, Theorem~4]{May1966},
with the odd-primary filtration of
\citet[Theorem~3.2.5]{Ravenel2004} and an additional filtration
adjustment constructed below.

A closer predecessor is the work of
\citet[author preprint, pp.~2--6]{LeeRavenel1994} on the nilpotence
order of the first beta-family sphere class $\beta_1$.
They use a modified May calculation on $V(3)$ and constrain
Adams--Novikov differentials through multiplication by $\beta_1$
and the $V(2)$-module structure of $V(3)$.
Our final obstruction calculation requires an annihilation relation
for products of a beta power with specific filtration-five classes,
together with a bound on the page at which those products vanish.
Toda's work on sphere relations \citep[p.~839, Corollary]{Toda1967}
and his subsequent cyclic extended-power constructions
\citep[Sections~1--3, especially Theorem~3]{Toda1968} provide the
topological antecedents. We give the particular coefficient-cohomology
calculations and finite-cell arguments needed for this bound below.

The construction has three stages. First we construct $Y=V(3)$ with a
binary product and a two-sided unit. These unit identities are enough
to turn a sphere map realizing $v_4$ into a self-map of $Y$. We obtain
the lower complexes and their products by finite cofiber constructions
and explicit vanishing of their attaching obstructions; we do not need
an associative product on $Y$.

Second, we prove that the filtration-zero class $v_4$ survives in the
Adams--Novikov spectral sequence for $Y$. To describe the obstruction
calculation, put $q=2p-2$ and $u_i=1+p+\cdots+p^{i-1}$ for $i\geq1$.
Divide the internal degree by $q$ and call the resulting integer the
\emph{weight}. A cobar class of cohomological degree $s$ and weight $w$
lies in stem $qw-s$. Every possible positive differential on $v_4$
therefore has target
\begin{equation}\label{intro:targets}
 (s,w)=(qk+1,u_4+k),\qquad k\geq1.
\end{equation}
A finite May calculation leaves only $k=1,2,p-1,p$. For $k=1$, a
relation supplied by the coboundary of $v_5$, together with an injective
multiplication on the target, forces the differential to vanish. For
$k=2,p-1$, we compute actual cohomology by taking the Frobenius kernel
of the group of truncated additive coordinate changes. The quotient
acts triangularly on kernel cohomology. Filtering this coefficient
module by weight identifies the first possible differential and
checks that both its source and its target are present. This step
controls the polynomial generators as well as the exterior generators
on the May page.

The remaining case $k=p$ requires a topological argument. We write
$\beta\in\pi_{pq-2}S^0_{(p)}$ for the first beta-family sphere class.
The entire $E_2$ target consists of products of a large power of its
detector with certain filtration-five classes of $Y$.
We realize those low-filtration classes by maps
$\gamma$ and construct an explicit relation
\[
 \beta^{\,p(p-2)+1}\gamma=0.
\]
The construction uses a cyclic extended power of a sphere and the
averaging retract of a smash power of a two-cell spectrum. The corresponding
Adams--Novikov product has filtration strictly smaller than the length of
the prospective
differential on $v_4$. Consequently the whole target has vanished
before that differential could occur. The bound on the page is the
point: eventual vanishing in homotopy, without such a bound, would
not settle this obstruction.

Finally the filtration-zero edge gives a map
$S^{2(p^4-1)}\to Y$ with Brown--Peterson Hurewicz image $v_4$.
Multiplication by the binary product on $Y$ extends it to
$\Sigma^{2(p^4-1)}Y\to Y$. Its cofiber is finite, and the
nonzerodivisor property of $v_4$ identifies its homology with the
required quotient. The bottom sphere map identifies the quotient
coaction and fixes the generator in degree zero.

Section~\ref{fw:section} fixes the Brown--Peterson coordinates and
the topological spectral-sequence conventions.
Section~\ref{may:preliminaries} constructs the filtered cobar calculation,
and Section~\ref{low:section} supplies the unital lower complexes.
Section~\ref{may:section} lists the height-four targets and removes the
first differential; Section~\ref{frob:section} proves the two Frobenius
vanishings. Section~\ref{top:section} establishes the homotopy relation,
its page bound, and the realizing cofiber. Appendix~\ref{app:resolution}
provides the ordinary convergent resolution and the filtered binary
pairing that connects the algebraic calculation to the actual
homotopy product.
